% GENERATED FILE. Edit canonical lessons, then run npm run generate:books.
% canonical-source-hash: fnv1a64:d104d421a98049f9
% canonical-lessons: SA-C166-asvatthah, SA-C166-nimbah, SA-C166-amravrksah, SA-C166-vatah, SA-C166-talah

\chapter{Sacred fig tree, Neem tree, Mango tree, Banyan tree, Palm tree}
\label{ch:sa-sacred-fig-tree-neem-tree-mango-tree-banyan-tree-palm-tree}
\hlchaptermodality{166}

\begin{chapteropening}
\textbf{By the end of this chapter:} \emph{I can say a sacred fig tree, a neem tree, a mango tree, a banyan tree and a palm tree.}

Complete the last lesson of chapter 166: I can say a sacred fig tree, a neem tree, a mango tree, a banyan tree and a palm tree.
\end{chapteropening}

\section[\texorpdfstring{\sk{अश्वत्थः} (aśvatthaḥ)}{अश्वत्थः (aśvatthaḥ)}]{\sk{अश्वत्थः} (aśvatthaḥ) — a sacred fig tree}

\label{lesson:SA-C166-asvatthah}



\begin{quote}
Before the new one: say the Sanskrit for a live coal, then the Sanskrit for iron.
\end{quote}

\subsection*{You'll want to know: \sk{अश्वत्थः}}
\textbf{\sk{अश्वत्थः}} — \emph{aśvatthaḥ} — \textquotedblleft{}a sacred fig tree\textquotedblright{}.

\begin{grammarlens}[title={What you've built}]
One more word for this chapter.
\end{grammarlens}

\subsection*{Your turn}
\begin{itemize}
\raggedright
  \item \emph{Say it:} \emph{aśvatthaḥ}
  \item \emph{Say it:} \emph{aśvatthaḥ}, once more
  \item \emph{Say it:} say \emph{aśvatthaḥ}
  \item \emph{Recall:} say the Sanskrit for a live coal, then the Sanskrit for iron, then say \emph{aśvatthaḥ} again
\end{itemize}

\begin{tcolorbox}[breakable,colback=teal!4,colframe=teal!35!black,title={Before you move on}]
What does \sk{अश्वत्थः} mean? (\textquotedblleft{}A sacred fig tree\textquotedblright{}.) Say it once more.
\end{tcolorbox}
\section[\texorpdfstring{\sk{निम्बः} (nimbaḥ)}{निम्बः (nimbaḥ)}]{\sk{निम्बः} (nimbaḥ) — a neem tree}

\label{lesson:SA-C166-nimbah}



\begin{quote}
Before the new one: say the Sanskrit for iron, then the Sanskrit for a sacred fig tree.
\end{quote}

\subsection*{You'll want to know: \sk{निम्बः}}
\textbf{\sk{निम्बः}} — \emph{nimbaḥ} — \textquotedblleft{}a neem tree\textquotedblright{}.

\begin{grammarlens}[title={What you've built}]
One more word for this chapter.
\end{grammarlens}

\subsection*{Your turn}
\begin{itemize}
\raggedright
  \item \emph{Say it:} \emph{nimbaḥ}
  \item \emph{Say it:} \emph{nimbaḥ}, once more
  \item \emph{Say it:} say \emph{nimbaḥ}
  \item \emph{Recall:} say the Sanskrit for iron, then the Sanskrit for a sacred fig tree, then say \emph{nimbaḥ} again
\end{itemize}

\begin{tcolorbox}[breakable,colback=teal!4,colframe=teal!35!black,title={Before you move on}]
What does \sk{निम्बः} mean? (\textquotedblleft{}A neem tree\textquotedblright{}.) Say it once more.
\end{tcolorbox}
\section[\texorpdfstring{\sk{आम्रवृक्षः} (āmravṛkṣaḥ)}{आम्रवृक्षः (āmravṛkṣaḥ)}]{\sk{आम्रवृक्षः} (āmravṛkṣaḥ) — a mango tree}

\label{lesson:SA-C166-amravrksah}



\begin{quote}
Before the new one: say the Sanskrit for a sacred fig tree, then the Sanskrit for a neem tree.
\end{quote}

\subsection*{You'll want to know: \sk{आम्रवृक्षः}}
\textbf{\sk{आम्रवृक्षः}} — \emph{āmravṛkṣaḥ} — \textquotedblleft{}a mango tree\textquotedblright{}.

\begin{grammarlens}[title={What you've built}]
One more word for this chapter.
\end{grammarlens}

\subsection*{Your turn}
\begin{itemize}
\raggedright
  \item \emph{Say it:} \emph{āmravṛkṣaḥ}
  \item \emph{Say it:} \emph{āmravṛkṣaḥ}, once more
  \item \emph{Say it:} say \emph{āmravṛkṣaḥ}
  \item \emph{Recall:} say the Sanskrit for a sacred fig tree, then the Sanskrit for a neem tree, then say \emph{āmravṛkṣaḥ} again
\end{itemize}

\begin{tcolorbox}[breakable,colback=teal!4,colframe=teal!35!black,title={Before you move on}]
What does \sk{आम्रवृक्षः} mean? (\textquotedblleft{}A mango tree\textquotedblright{}.) Say it once more.
\end{tcolorbox}
\section[\texorpdfstring{\sk{वटः} (vaṭaḥ)}{वटः (vaṭaḥ)}]{\sk{वटः} (vaṭaḥ) — a banyan tree}

\label{lesson:SA-C166-vatah}



\begin{quote}
Before the new one: say the Sanskrit for a neem tree, then the Sanskrit for a mango tree.
\end{quote}

\subsection*{You'll want to know: \sk{वटः}}
\textbf{\sk{वटः}} — \emph{vaṭaḥ} — \textquotedblleft{}a banyan tree\textquotedblright{}.

\begin{grammarlens}[title={What you've built}]
One more word for this chapter.
\end{grammarlens}

\subsection*{Your turn}
\begin{itemize}
\raggedright
  \item \emph{Say it:} \emph{vaṭaḥ}
  \item \emph{Say it:} \emph{vaṭaḥ}, once more
  \item \emph{Say it:} say \emph{vaṭaḥ}
  \item \emph{Recall:} say the Sanskrit for a neem tree, then the Sanskrit for a mango tree, then say \emph{vaṭaḥ} again
\end{itemize}

\begin{tcolorbox}[breakable,colback=teal!4,colframe=teal!35!black,title={Before you move on}]
What does \sk{वटः} mean? (\textquotedblleft{}A banyan tree\textquotedblright{}.) Say it once more.
\end{tcolorbox}
\section[\texorpdfstring{\sk{तालः} (tālaḥ)}{तालः (tālaḥ)}]{\sk{तालः} (tālaḥ) — a palm tree}

\label{lesson:SA-C166-talah}



\begin{quote}
Before the new one: say the Sanskrit for a mango tree, then the Sanskrit for a banyan tree.
\end{quote}

\subsection*{You'll want to know: \sk{तालः}}
\textbf{\sk{तालः}} — \emph{tālaḥ} — \textquotedblleft{}a palm tree\textquotedblright{}.

\begin{grammarlens}[title={What you've built}]
One more word for this chapter.
\end{grammarlens}

\subsection*{Your turn}
\begin{itemize}
\raggedright
  \item \emph{Say it:} \emph{tālaḥ}
  \item \emph{Say it:} \emph{tālaḥ}, once more
  \item \emph{Say it:} say \emph{tālaḥ}
  \item \emph{Recall:} say the Sanskrit for a mango tree, then the Sanskrit for a banyan tree, then say \emph{tālaḥ} again
\end{itemize}

\begin{tcolorbox}[breakable,colback=teal!4,colframe=teal!35!black,title={Before you move on}]
What does \sk{तालः} mean? (\textquotedblleft{}A palm tree\textquotedblright{}.) Say it once more.
\end{tcolorbox}
